%% =====================================================================
%%  1 Introduction  (median 1.5 pages; 70 of 74 gold papers use this title)
%% =====================================================================
%%  GUIDE
%%   - Para 1: the physical situation and why it matters. Two or three
%%     sentences, no history of physics.
%%   - Para 2: what is already known, with citations [1-3]. Name the
%%     specific assumption of each prior model, not just "many studies".
%%   - Para 3: the gap this problem exposes, and your contribution in one
%%     sentence: "In this paper we derive ..., validate it against ...,
%%     and obtain ...."
%%   - Close with a two-sentence roadmap, not a paragraph.
%%  Do not write "since the beginning of time" or a popular-science intro.
%% =====================================================================
\section{Introduction}
\label{sec:intro}

A breach in the pressure hull of a crewed module vents its atmosphere in
minutes, and the resulting pressure--time history decides whether the crew
can seal the hole, reach a shelter, or must abandon the module. The design
question is therefore quantitative: given the hole size and the module
volume, how long does the crew have before the pressure falls below the
survival threshold, and how much does that time depend on the parameters
that are known only roughly?

Compressible-flow theory gives the mass flux through an orifice, in both
the choked and the subsonic regime \citep{shapiro1953,anderson2003}, and
the quasi-steady reservoir model built on it is standard in vacuum
technology \citep{landau1987}. Applied to a habitat, that standard
treatment usually stops at two simplifications: the discharge coefficient
is taken from a handbook, and the gas is assumed isothermal or adiabatic
throughout. Neither is harmless here. The handbook value belongs to a
sharp-edged plate, not to a hole in a hull, and the thermal state of the
gas controls its density, hence the mass that still has to leave and the
time the crew actually has.

In this paper we (i) derive a quasi-steady blow-down model that keeps the
choked branch and adds wall heat transfer as a first-order relaxation,
(ii) calibrate its single free parameter, the discharge coefficient,
against published leak-down data, and (iii) reduce the result to a scaling
law, $t_{0.3}=1.41\,\tau$ with $\tau\propto d^{-2}$, that a crew can
evaluate without a computer. Section~\ref{sec:model} sets out the model,
Section~\ref{sec:results} answers the three questions of the problem
statement, and Section~\ref{sec:discussion} quantifies the uncertainty and
states where the model must not be used.
